\section{生产级韧性与安全重构：设计说明与面试答法}

\begin{summaryBox}{这一章回答什么}
这一章不是把安全名词罗列一遍，而是把项目从“功能能够跑”提升到“外部依赖失败时仍可解释、数据越权时默认拒绝、状态恢复时有事实依据”的设计完整讲清楚。每个主题都按同一套问题展开：为什么需要它，采用什么方案，为什么没有选其他方案，收益是什么，代价是什么，已经考虑了哪些边界，还有哪些边界不能声称已经解决。
\end{summaryBox}

\begin{riskBox}{面试中的总原则}
当前项目是单用户自托管系统，但单用户不等于无认证，也不等于可以省略 owner scope。面试时应把“已实现”“已测试”“设计预留”和“尚未完成”分开说。尤其是生产密钥、HTTPS、真实外部模型和多主体部署，不能用本地单元测试替代现场证据。
\end{riskBox}

\subsection{可退化架构：从全有或全无变成健康、降级、恢复}

\begin{qaBox}{面试官可能问：为什么不让依赖失败就直接报错？}
因为研究系统的能力不是同质的。PostgreSQL、认证密钥和审计写入决定安全事实，失败时必须拒绝服务；RAG、重排序、浏览器和备用模型则是可选能力，失败时可以缩小研究范围并继续给出明确标注的结果。把所有依赖放进一个 try/except 会把安全依赖和可选能力混在一起，既可能扩大安全风险，也会把一次局部故障放大成整个应用不可用。
\end{qaBox}

\paragraph{设计。}
在 \term{app/security/capabilities.py} 中把外部能力抽象成 \term{CapabilityHealth} 状态机：\term{healthy} 表示最近探测成功，\term{degraded} 表示最近调用失败但应用有替代路径，\term{recovering} 表示冷却期后正在进行单飞探测。首次使用才初始化能力；失败记录 \term{reason\_code} 与时间；默认冷却 30 秒；并发请求共享一次恢复探测，避免所有请求同时下载模型或连接死服务。

降级结果不能只返回空列表，而应携带 \term{capability}、\term{state}、\term{reason\_code}、\term{retry\_after\_seconds} 和 \term{fallback\_used}。因此 DAG、SSE、审计和最终报告都能回答三个问题：哪个能力没用上，为什么没用上，什么时候可以再试。RAG 失败时跳过内部检索并标注“未使用内部知识库”；reranker 失败时保留向量/BM25 结果；浏览器失败时仍可保留搜索；报告流在首个正文输出前才允许切备用模型。

\paragraph{备选方案与取舍。}
方案 A 是任一依赖失败即终止。它的优点是语义简单、不会产生“部分结果”，缺点是可用性差、局部故障会放大，而且用户不知道已有证据是否仍然有价值。方案 B 是每个请求临时判断依赖是否可用。它比全局终止灵活，但并发下容易产生探测风暴，且不同节点可能对同一个能力做出不一致判断。方案 C 是集中式能力状态机加按能力定义的降级矩阵。项目选择方案 C，因为它同时解决状态一致性、节流恢复和用户可解释性；代价是增加状态维护、测试矩阵和错误码设计。

\paragraph{收益、代价与边界。}
收益是故障半径可控、恢复不需要重启、SSE 能继续表达进度，且安全关键依赖不会为了可用性而放行。代价是“部分完成”比成功/失败二元状态更难测试，也要求报告诚实标注证据缺口。已经考虑了冷却、单飞、可选能力回退和最终报告边界；尚未完成的是生产级指标告警、跨进程共享能力状态和真正外部模型的现场恢复演练。当前测试证明的是本进程状态机和降级语义，不是第三方服务的 SLO。

\subsection{安全边界：数据归属、令牌生命周期与默认拒绝}

\begin{qaBox}{面试官可能问：单用户系统为什么还要 owner scope？}
因为单用户是部署规模，不是数据访问模型。只要存在研究会话、文档、技能、配置、审批和审计，就需要明确它们属于哪个主体。严格的 owner 条件可以阻止历史数据被错误读取，也让未来增加第二主体时不必重写所有查询。
\end{qaBox}

\paragraph{设计。}
认证子系统使用 bcrypt 保存密码哈希，首次初始化通过 PostgreSQL advisory transaction lock 串行化，保留用户名拒绝被用于伪装系统主体。服务端 session 和 API token 只保存哈希值；明文 token 仅在创建响应中返回一次。每次校验同时检查撤销时间、过期时间、主体启用状态和 \term{session\_version}；密码轮换或统一撤销时递增版本并撤销旧 session/token，因此旧凭据不能继续访问。

所有资源查询使用主体条件，跨主体对象统一返回 404，避免泄露对象存在性。认证、审批和审计属于关键安全状态，不能因为 Redis 缓存故障而退化为放行。安全中间件统一设置精确 CORS、CSP、点击劫持防护、Referrer Policy、Content-Type 防嗅探和安全 Cookie；认证、上传和 token 路由还应使用独立限流桶。

\paragraph{备选方案与取舍。}
方案 A 只绑定 localhost，部署简单但无法安全远程访问。方案 B 完全信任反向代理，性能好但应用自身失去独立安全边界，代理配置错误会直接变成越权。方案 C 使用应用层密码、服务端 session 和可撤销 API token，并把 HTTPS 代理作为额外防线。项目采用方案 C，因为它可独立部署、能撤销泄露凭据，且不把安全责任转嫁给代理；代价是要维护认证表、锁、轮换和限流逻辑。

\paragraph{已经考虑与尚未考虑。}
已经考虑了 session fixation、密码轮换、删除/停用主体后的旧 token 失效、跨用户资源 404、Redis 故障不放行和 owner scope 回归。尚未在本轮引入多用户角色、组织/租户层级、WebAuthn、设备绑定和集中身份提供商；这些不能在面试中说成当前功能。

\subsection{配置与密钥：按主体隔离，并讲清威胁模型}

\begin{qaBox}{面试官可能问：应用层加密到底防什么？}
它主要防数据库文件、备份和日志泄露后直接暴露 LLM API Key。它不防已经获得进程权限的攻击者，因为运行中的进程必须能够解密并调用模型。把保护对象和不保护对象说清楚，比声称“密钥绝对安全”更可信。
\end{qaBox}

\paragraph{设计。}
LLM 配置属于受认证保护的管理员资源，数据库中的 API Key 使用部署主密钥加密，并以 \term{enc:} 前缀区分已迁移值。读取和写入均是幂等的，旧明文可以渐进迁移，不会因为重复启动而重复加密。配置 key 使用 \term{llm:<owner\_id>}，查询同时带 \term{owner\_id}，旧单管理员 \term{llm} 记录只在同一 owner 范围内兼容读取。生产模式拒绝缺失或示例主密钥、不安全 Cookie 和通配 CORS。

\paragraph{方案比较。}
明文入库最简单但数据库备份即等于泄露；只放环境变量能减少数据库暴露面，却失去运行时配置和管理员资源审计；应用层加密兼顾运行时配置、备份安全和渐进迁移，但引入主密钥轮换、启动失败和密钥丢失恢复问题。项目选择第三种，并明确“主密钥不与数据库备份放在一起”。

\paragraph{收益、代价与遗漏。}
收益是泄露数据库不直接泄露模型密钥，配置更新可审计且不必重训练。代价是主密钥轮换需要版本化密文和迁移工具，主密钥丢失会让密文不可恢复。已经考虑了前缀幂等迁移、数据库/备份威胁模型和生产启动检查；尚未完成 KMS/HSM 托管、自动轮换、密钥版本回滚和真实生产密钥演练。

\subsection{兼容层与 facade：重构不应一次性破坏调用面}

\term{app/agent\_tools} 的兼容层保留旧模块导入面，再逐步把解析、schema、执行和具体工具实现拆到更清晰的边界。它解决的是重构风险，不是业务功能本身：旧调用方继续拿到稳定 facade，新代码可以使用更细粒度的 adapter、policy 和执行器，迁移完成后再逐步收窄旧接口。

方案 A 是一次性重写所有调用方，结构最干净，但变更面大、回归定位困难；方案 B 是永远保留一层万能兼容代码，迁移成本低，却会形成双重语义和死代码；方案 C 是带弃用策略、契约测试和观测的渐进 facade。项目采用 C。收益是降低重构破坏性、支持分阶段发布和回滚；代价是短期内需要维护两套入口、增加文档和测试成本。已经考虑了旧导入兼容和 schema/执行边界；尚未完成的是统一弃用时间表、调用方迁移完成度指标和删除旧 facade 的最终版本。

\subsection{统一 LLM adapter：把 provider 差异隔离在边界}

\paragraph{设计。}
业务 Agent 不直接判断 provider，而通过 endpoint resolver 与 adapter 获取规范化的 base URL、认证头、模型名和客户端。OpenAI 兼容端点、Ollama 与 Anthropic 的差异集中在 \term{app/endpoint\_resolver.py} 和 \term{app/llm\_adapters.py}。HTTP 客户端共享连接池，并分别配置 connect/read/write/pool timeout；死端点进入冷却，有限重试只对可恢复的网络错误生效。

非流式调用可以在首个响应前切备用模型；流式调用把“是否已经向用户输出正文”作为协议状态。一旦已有正文，主模型失败只能发出可恢复的终止错误，不能无提示地切到备用模型，否则用户会看到语义断裂的半份报告。usage 缺失时使用保守估算，并在审计中标记 \term{estimated}。

\paragraph{为什么不是业务层 if/else。}
把 provider 判断散落在 Agent 中，短期代码少，长期会让 URL、认证头、超时、模型过滤和 fallback 规则彼此漂移。统一 adapter 的优点是替换模型和增加 provider 的影响面小、协议测试集中；缺点是 adapter 本身成为高复杂度边界，错误归类和原生 provider 特性可能被过度抽象。当前已覆盖 URL、认证头、模型筛选、usage 估算和首字节规则；尚未覆盖所有供应商的真实流式协议、工具调用差异和跨区域熔断指标。

\subsection{提示注入：把不可信数据放进数据流，而不是一句 system prompt}

网页、邮件、上传文档、RAG 片段、Skill 外部文本以及 MCP/tool 输出都必须携带来源、信任级别和内容边界。它们只能作为明确标注的用户上下文，不能拼接进 system role，也不能把其中的“请执行某命令”解释为控制指令。\term{prompt\_security} 负责数据封装，guardrail 负责研究政策，两者要分开建模：前者解决“这段文本能不能影响指令解释”，后者解决“这个问题是否允许联网、写入或输出”。

收益是攻击者无法仅靠网页文本升级为系统指令，且审计可以追踪污染来源；代价是 prompt 更长、模型可能降低可读性，且标记本身不是形式化证明。已经考虑了常见外部文本边界和 system role 隔离；尚未完成跨模型的注入基准、工具参数级 taint tracking、图像/多模态载荷和红队样本库。

\subsection{健康检查：liveness 与 readiness 必须分开}

\term{/health} 只回答进程是否活着，不能因为数据库暂时不可用就让容器被重启风暴拖垮。\term{/ready} 回答“能不能接流量”：数据库可连、schema 可用、认证主密钥有效、审计可写时才返回 200；RAG、Redis、浏览器和备用模型以不泄露敏感细节的信息状态返回。Docker healthcheck 应指向 \term{/ready}，因为编排系统需要依据接流量能力摘除实例。

方案 A 只有一个 health endpoint，简单但无法区分进程存活与业务可用；方案 B 只检查 HTTP server，误把“能响应”当成“能安全服务”；方案 C 分离 liveness/readiness 并分级关键依赖。项目采用 C，代价是部署和测试需要理解两种状态，且 readiness 不能泄露连接串、密钥或内部异常堆栈。当前代码已覆盖数据库、Redis 和认证状态的基本展示；真实容器冷启动和代理转发仍需现场验证。

\subsection{上传与文档生命周期：不是一个 save(file)}

上传链路应按顺序执行认证、owner 校验、文件名与路径规范化、扩展名/MIME/魔数交叉校验、大小/页数/解压限制、文本提取超时、哈希去重、主体配额和按主体限流。成功入库后，source 与 chunks 使用同一 owner；删除 source 时同步删除其 chunks 和向量记录；检索查询默认带 owner 条件。不要把用户原始路径写入宿主机，也不要用文件扩展名代替内容检测。

收益是降低路径逃逸、伪造 MIME、压缩炸弹、重复存储和跨主体检索风险；代价是解析库、缩略图、EXIF 处理和异步清理会增加复杂度，严格限制也可能拒绝合法的大文件。已经考虑了路径、MIME、哈希、owner、限流和同步删除；尚未完成的重点是对每一种文档格式做系统 fuzzing、病毒扫描、对象存储生命周期和可恢复的异步清理队列。

\subsection{备份导出/导入：merge-only 仍然要有事务边界}

导出只包含当前管理员拥有的认可资源，不导出密钥、session、token 和审计秘密。导入先做版本号与白名单结构校验，再按内容哈希去重、补齐 owner，并在同一个数据库连接的单一 transaction 中完成整批写入。任一 source 失败，整个批次回滚；成功响应同时返回 imported 与 skipped，保证幂等和冲突可见。

覆盖写入看起来简单，但会丢失当前数据、破坏 owner 关系，失败中途还很难恢复；逐条提交虽然易实现，却会留下半批导入。merge-only + 单事务的代价是大批次长事务占用连接、锁和 WAL，需要进一步做批量上限、分段导入协议或 staging table。当前已实现版本化、白名单、哈希 merge、owner 绑定和整批回滚测试；尚未完成跨版本迁移、加密备份容器、断点续传和大规模导入压测。

\subsection{本地持久化：JSON 也要遵守数据库式原子写}

仍需写本地文件的配置、Harness 状态和 session 快照，应统一使用同目录临时文件、完整写入、\term{flush}、\term{fsync} 和 \term{os.replace}。同目录保证替换发生在同一文件系统，\term{os.replace} 保证读者不会看到半截 JSON。写入工具应统一处理权限、编码、临时文件清理和 Windows 行为，禁止模块各自随手 \term{json.dump()}。

原子写不能解决磁盘损坏、错误数据逻辑或跨文件事务；它只保证单文件替换边界。因此 Harness 与数据库双写仍需分域 source of truth 和补偿恢复。收益是断电、进程终止和并发读时不容易得到损坏文件；代价是多一次磁盘写入、fsync 延迟和磁盘空间峰值。当前应把这套规则作为所有新增本地持久化的硬约束，而不能声称它等价于完整数据库事务。

\subsection{测试文化：每个真实事故都留下可回归的名字}

测试不应只验证 happy path，而应围绕曾经发生或可能发生的故障命名：初始化并发、错误密码、旧 token 失效、跨主体读取、短暂依赖故障、30 秒冷却、单飞恢复、首字节后禁止切换、恶意上传、伪造 MIME、导入回滚、提示注入和 readiness 拒流量。测试名本身就是维护文档，告诉未来维护者为什么一个看似多余的分支必须保留。

覆盖率只是信号，不是安全证明。90\% 以上可以降低未执行路径风险，但不能证明 threat model 完整；必须结合 property-based/fuzz 测试、迁移升级测试、真实协议测试和红队样本。当前自动化验证覆盖安全/adapter 94\%、全仓测试 172 项通过；尚未完成的是生产依赖现场、容器冷启动、多进程并发和长期故障注入。

\subsection{不在本轮范围的后台自动化与未来多主体}

定时任务、事件总线、邮件、CalDAV、跨设备同步和复杂角色体系不属于本轮生产重构边界。面试时可以说明其设计借鉴点：调度器要处理时区、短月钳制、幂等、单飞、串行执行和重启触发保留；事件必须按 owner 路由，不能一个用户的事件触发所有人的任务。但当前项目没有把这些能力伪装成已完成功能。保留 owner schema 和审计边界，是为未来扩展提供兼容接口，而不是提前承诺多用户产品。

\subsection{最终取舍总表}

\begin{center}
\small
\begin{tabularx}{\textwidth}{>{\bfseries}lXXX}
\toprule
主题 & 采用方案 & 主要收益 & 主要未覆盖项 \\
\midrule
依赖故障 & 能力状态机 + 分级降级 & 可用性与安全边界分离 & 跨进程状态、真实 SLO \\
访问控制 & 应用层认证 + owner scope & 可撤销、可独立部署 & 多角色、WebAuthn \\
密钥 & 应用层加密 + 渐进迁移 & 降低数据库泄露影响 & KMS/HSM、自动轮换 \\
LLM & 统一 adapter + 首字节协议 & provider 可替换、流式不串语义 & 全 provider 现场协议 \\
上传 & 生命周期校验 + 哈希 + owner & 减少文件攻击和越权 & fuzz、病毒扫描、对象存储 \\
备份 & 白名单 merge + 单批事务 & 幂等且不破坏现有数据 & 大规模断点续传 \\
健康 & liveness/readiness 分离 & 编排语义准确 & 生产代理现场演练 \\
测试 & 事故驱动回归 + 覆盖率 & 故障可复现、边界可审计 & 长期故障注入 \\
\bottomrule
\end{tabularx}
\end{center}

\begin{qaBox}{两分钟收口答案}
我没有把所有外部依赖都当成同一种错误，而是先按安全关键程度分级：数据库、认证密钥和审计写入失败时拒绝服务，RAG、浏览器、重排序和模型端点失败时通过能力状态机进入降级并持续探测恢复。所有资源查询坚持 owner scope，token 只存哈希且支持版本化撤销；API Key 用部署主密钥加密，但我会明确它防的是数据库/备份泄露，不防已攻陷进程。LLM provider 差异集中在 adapter，流式响应一旦输出正文就不允许静默切模型。上传、备份、提示注入和健康检查都有明确边界。最后我会把已测试能力、设计预留和尚未完成的生产现场演练分开说，而不是用覆盖率或文档替代真实验证。
\end{qaBox}
